% Clustering and the transfer-operator gap — machine-checked equivalence paper.
% Lane: O (operator bridge), modules YangMills/OS/**.
% REGIME: tricotomy; NO scores, NO desk language, NO commit chronology beyond
% the reproducibility freeze.  Freeze: Lean tree of source checkpoint
% ad9c93d793fa0b082f2373c2aea3ea5b99f9c6b8, full core build 8415 jobs,
% oracle transcripts named by their code commits (see Section 7).
\documentclass[11pt]{article}
\usepackage[margin=1.1in]{geometry}
\usepackage{amsmath,amssymb,amsthm}
\usepackage{booktabs,longtable,array}
\usepackage[colorlinks=true,linkcolor=blue,urlcolor=blue,citecolor=blue]{hyperref}

\newtheorem{theorem}{Theorem}[section]
\newtheorem{lemma}[theorem]{Lemma}
\newtheorem{proposition}[theorem]{Proposition}
\newtheorem{corollary}[theorem]{Corollary}
\newtheorem{remark}[theorem]{Remark}
\theoremstyle{definition}
\newtheorem{definition}[theorem]{Definition}

\emergencystretch=3em
\tolerance=2000

\newcommand{\lean}[1]{\texttt{#1}}
\newcommand{\anchor}{ad9c93d793fa0b082f2373c2aea3ea5b99f9c6b8}
\newcommand{\osfile}[2]{%
  \href{https://github.com/lluiseriksson/THE-ERIKSSON-PROGRAMME/blob/\anchor/YangMills/OS/#1.lean}{\nolinkurl{#2}}}
\newcommand{\repofile}[2]{%
  \href{https://github.com/lluiseriksson/THE-ERIKSSON-PROGRAMME/blob/\anchor/#1}{\nolinkurl{#2}}}
\newcommand{\R}{\mathbb{R}}
\newcommand{\C}{\mathbb{C}}
\newcommand{\SU}{\mathrm{SU}}
\newcommand{\Om}{\Omega}
\newcommand{\inn}[2]{\langle #1,\, #2\rangle}
\newcommand{\ket}[1]{|#1\rangle}
\newcommand{\bra}[1]{\langle #1|}

\title{Clustering and the Transfer-Operator Gap:\\
A Machine-Checked Dense-Family Criterion\\
\large Observable-dependent prefactors are harmless\\
at a common exponential rate}
\author{Lluis Eriksson}
\date{July 27, 2026 (version 1.2)}

\begin{document}
\maketitle

\begin{abstract}
Inside a Lean~4 formalization programme for four-dimensional $\SU(N_c)$
lattice Yang--Mills, we machine-check the operator-theoretic criterion
that stands between exponential decay of a Euclidean correlator and a
spectral gap of a transfer operator.  Let $T$ be a bounded self-adjoint
operator on a Hilbert space and $\Om$ a unit vector \emph{fixed by} $T$,
so that $T\Om = \Om$, and put $S = T - \ket{\Om}\bra{\Om}$.  Exponential
decay at rate $r$ of the \emph{connected} two-point function
$\inn{v}{T^n v} - |\inn{\Om}{v}|^2$ at every $v$ is \emph{equivalent} to
the operator-norm bound $\|S\| \le r$.

The substantive part is the dense-family criterion.  Writing
$\mathcal{D}_r = \{v : \exists C,\ \forall n,\ \|S^n v\| \le C r^n\}$, we
prove that $\mathcal{D}_r$ is a linear subspace and that its density
alone forces $\|S\| \le r$, the constants being entirely unconstrained:
a family of observables whose span is dense, each carrying its own finite
constant, suffices.  Consequently prefactors that grow with the support
of the observable --- the shape cluster expansions produce --- do
\emph{not} obstruct the gap, provided the exponential \emph{rate} is
common to the family and the family spans densely.  Those two provisos
are essential; without them the statement is false.

No mathematical novelty is claimed for the criterion itself, which we
expect to be known in the language of local spectral theory; what is
offered is its mechanization, its packaging for families of observables,
and the consequence for prefactors.  We also record what the
formalization does \emph{not} contain: no Osterwalder--Seiler Hilbert
space for any gauge theory, no reflection positivity of the Wilson
measure, no identification of a Euclidean correlator with a matrix
element.  Nothing here is a claim about the continuum limit or about the
Clay problem.  All results are machine-checked with no \lean{sorry} and
no project axioms.
\end{abstract}

\section{What is proved, and what is not}

\subsection*{The gap in the chain}

A lattice mass gap, in the sense of Osterwalder and
Seiler~\cite{OsterwalderSeiler}, is a statement about the spectrum of an
operator: the transfer operator $T$ obtained from the Euclidean measure
by reflection positivity satisfies $0 \le T \le 1$, fixes the vacuum
vector $\Om$, and has a spectral gap on $\Om^\perp$.  Exponential decay
of a Euclidean correlator is, by itself, a statement about a
real-valued function of the separation.  The two are connected by a
theorem, not by terminology.

Within the formalization programme this paper belongs to, that theorem
was absent.  The terminal statement of the lattice assembly concludes
about a function $\mathrm{cov} : \mathbb{N} \to \R$; the programme
contains no Hilbert space of states, no transfer operator, no reflection
positivity, and no spectrum.  This paper supplies the \emph{autonomous}
half of the missing link: the operator-theoretic equivalence, at the
level of generality where it is a theorem, machine-checked, together
with the sharpening that determines what an application must supply.

\subsection*{Scope, stated in advance}

The following are \emph{not} established in this paper, nor anywhere in
the Lean development accompanying it:
reflection positivity of the Wilson gauge measure; the
Gelfand--Naimark--Segal quotient; the existence of a transfer operator
for any gauge theory; the identification
$\mathbb{E}[A \cdot \theta_n B] = \inn{A\Om}{T^n B\Om}$; density of any
concrete family of gauge-invariant observables in the fluctuation sector
of any such space.  Consequently nothing in this paper is a claim about
Yang--Mills, about the continuum limit, or about the Clay problem, and no
statement below should be read as transferring importance from the
subject matter of the programme to the theorems proved here.  The
equivalence itself is classical; see Section~\ref{sec:prov}.

\section{Setting}

Let $H$ be a Hilbert space over $\R$ (Sections~\ref{sec:equiv}) or over
$\C$ (Sections~\ref{sec:sharp}--\ref{sec:vu}).

\begin{definition}[transfer data]
\label{def:data}
A pair $(T,\Om)$ with $T$ a bounded operator on $H$ and $\Om \in H$ is
\emph{transfer data} when $T$ is symmetric, $\|\Om\| = 1$, and
$T\Om = \Om$.  We write $P = \ket{\Om}\bra{\Om}$ for the rank-one
projection $Pv = \inn{\Om}{v}\,\Om$ and
\[
  S \;=\; T - P
\]
for the \emph{vacuum-projected} operator.  These are the properties that
reflection positivity is used to produce; here they are hypotheses.
\end{definition}

\begin{definition}[connected two-point function]
\label{def:conn}
For $v \in H$ and $n \in \mathbb{N}$,
\[
  G_v(n) \;=\; \inn{v}{T^n v} \;-\; \inn{v}{\Om}\inn{\Om}{v}.
\]
\end{definition}

The elementary reason the \emph{truncated} correlator, and not the full
one, is the object that sees a gap is the following identity, which is
proved by induction from $S\Om = 0$ and $\inn{\Om}{Sv} = 0$.

\begin{lemma}[structural identity]
\label{lem:struct}
For transfer data and every $n \ge 0$,
$S^{n+1} v = T^{n+1} v - \inn{\Om}{v}\,\Om$, and hence
$G_v(n+1) = \inn{v}{S^{n+1} v}$.
\end{lemma}

\section{The equivalence}
\label{sec:equiv}

\begin{theorem}[clustering $\iff$ gap]
\label{thm:equiv}
Let $(T,\Om)$ be transfer data on a real Hilbert space and $r \ge 0$.
Then
\[
  \|S\| \;\le\; r
  \qquad\Longleftrightarrow\qquad
  \forall v \in H,\ \exists C,\ \forall n,\quad |G_v(n)| \;\le\; C\,r^n .
\]
\end{theorem}

Two features of the statement are deliberate.  First, the conclusion is
an operator-norm bound and not an assertion of the form
$\exists\,m > 0$; the latter shape admits vacuous witnesses, and this
programme's own history contains one.  Second, the statement is an
equivalence, so it cannot be weakened into vacuity in either direction.

\begin{proof}[Proof sketch]
\emph{Forward.}  For $n \ge 1$, Lemma~\ref{lem:struct} gives
$|G_v(n)| = |\inn{v}{S^n v}| \le \|v\|\,\|S^n\|\,\|v\| \le \|v\|^2 r^n$;
the $n = 0$ term is $\|v\|^2 - |\inn{\Om}{v}|^2 \in [0,\|v\|^2]$ by
Cauchy--Schwarz.

\emph{Backward, the case $r = 0$.}  Lemma~\ref{lem:struct} at $n = 1$
gives $0 \le \|Sv\|^2 = G_v(2) \le C_v\cdot 0^2 = 0$ for every $v$, so
$S = 0$ and $\|S\| \le 0$.

\emph{Backward, the case $r > 0$.}  For $n \ge 1$ --- and only for
$n \ge 1$, since $G_v(0) \ne \inn{v}{S^0 v}$ in general ---
Lemma~\ref{lem:struct} and self-adjointness give
$G_v(2n) = \inn{v}{S^{2n}v} = \|S^n v\|^2$, so clustering at even times
gives $\|S^n v\| \le \sqrt{C_v}\,r^n$ pointwise.  Banach--Steinhaus
applied to $\{r^{-n}S^n : n \ge 1\}$ upgrades this to $\|S^n\| \le M r^n$;
and for symmetric $S$ one has $\|S^2\| = \|S\|^2$ --- proved from
Cauchy--Schwarz, without the $C^*$-algebra API --- hence
$\|S\|^{2^k} \le M r^{2^k}$ for all $k$, and $\|S\| \le r$.
\end{proof}

\section{The sharp form: the constants do not matter}
\label{sec:sharp}

Theorem~\ref{thm:equiv} assumes clustering at \emph{every} vector.  A
cluster expansion never delivers that: it bounds a \emph{family} of
observables.  On a merely dense set, the elementary argument --- extend
the quadratic-form bound by continuity --- requires the constant to be
uniformly quadratic, $C_v = K\|v\|^2$.  That requirement is a property of
that proof.  It is not necessary.

\begin{definition}[decay domain]
\label{def:decay}
For a bounded operator $S$ on $H$ and $r \ge 0$,
\[
  \mathcal{D}_r(S) \;=\; \bigl\{\, v \in H \;:\; \exists C \ge 0,\
  \forall n,\ \|S^n v\| \le C\,r^n \,\bigr\}.
\]
(The Lean definition leaves $r$ and $C$ unconstrained in $\R$ and carries
$r \ge 0$ as a hypothesis at each use site, which is why the subspace
lemma below holds for either sign; the restriction here is for
readability.)
\end{definition}

\begin{lemma}
\label{lem:submodule}
$\mathcal{D}_r(S)$ is a linear subspace of $H$ (a \lean{Submodule} in the
formalization).
\end{lemma}

\begin{proof}
$C_{u+v} \le C_u + C_v$ and $C_{\lambda v} = |\lambda|\,C_v$.
\end{proof}

Lemma~\ref{lem:submodule} contains the algebraic step that passes from a
controlled family to its linear span; the substantive mathematics is the
functional-calculus argument of Theorem~\ref{thm:sharp}.  What the lemma
buys is that
per-observable constants, however fast they grow along any exhaustion of
the space, close up under linear combination with no uniformity
whatsoever.

\begin{theorem}[sharp form]
\label{thm:sharp}
Let $S$ be a self-adjoint bounded operator on a complex Hilbert space
$H \neq 0$ and $r \ge 0$.  If $\mathcal{D}_r(S)$ is dense in $H$, then
$\|S\| \le r$.  Consequently, if $D \subseteq H$ has dense span and every
$v \in D$ admits some finite $C_v$ with $\|S^n v\| \le C_v r^n$ for all
$n$, then $\|S\| \le r$.
\end{theorem}

\begin{proof}
Suppose $\|S\| > r$.  The norm is attained on the spectrum, so there is
$\lambda_0 \in \sigma(S)$ with $|\lambda_0| = \|S\|$; choose
$r < c < |\lambda_0|$ and set $f(x) = \max(0, |x| - c)$, so that $f$
vanishes on $\{|x| \le c\}$ and $f(\lambda_0) > 0$.  Define
\[
  k_n(x) \;=\; \frac{f(x)}{(\max(c,|x|))^{2n}} .
\]
The denominator is $\ge c^{2n} > 0$ for \emph{every} real $x$, so $k_n$
is continuous with no case analysis and no division by zero, even when
$0 \in \sigma(S)$; and $k_n(x)\,x^{2n} = f(x)$ identically, the even
power letting $|x|$ stand in for $x$.  Hence, in the continuous
functional calculus, $f(S) = k_n(S)\,S^{2n}$, and on the spectrum
$\|k_n\|_\infty \le \|S\|/c^{2n}$.  For $v \in \mathcal{D}_r(S)$,
\[
  \|f(S)v\| \;=\; \|k_n(S)\,S^{2n}v\|
  \;\le\; \frac{\|S\|}{c^{2n}}\,C_v\,r^{2n}
  \;=\; \|S\|\,C_v \Bigl(\frac{r}{c}\Bigr)^{2n}
  \;\xrightarrow[n\to\infty]{}\; 0 ,
\]
so $f(S)$ annihilates $\mathcal{D}_r(S)$, hence all of $H$ by density and
continuity.  But $\|f(S)\| \ge |f(\lambda_0)| > 0$, a contradiction.  The
second statement follows from Lemma~\ref{lem:submodule}, which places the
span of $D$ inside $\mathcal{D}_r(S)$.
\end{proof}

\subsection*{Why a functional calculus is needed}
\label{rem:notelem}

Three routes that suggest themselves all fail.  Recording why is the main
reason a formalization of this criterion is worth reading: it says which
architecture survives, not merely that a theorem type-checks.

\smallskip\noindent\emph{Banach--Steinhaus does not apply.}  A dense
\emph{subspace} may be meagre --- the span of a countable family is --- so
the non-meagreness hypothesis is not at hand, and pointwise bounds on such
a set yield no uniform bound.  This is exactly the difference between
Theorem~\ref{thm:equiv}, where the hypothesis is at every vector of $H$
and completeness pays for the uniformity, and Theorem~\ref{thm:sharp},
where it is not.

\smallskip\noindent\emph{Powers do not suffice.}  Put
$B = S^2/\|S\|^2$, a positive contraction with $\|B\| = 1$.  It can happen
that $B^n v \to 0$ for \emph{every} $v$ while $\|B^n\| = 1$ for every $n$
--- multiplication by $x$ on $L^2[0,1]$ is exactly that --- so the norm
level alone yields no contradiction.  (For a general self-adjoint $S$ the
powers need not tend to $0$ at all: if $\pm\|S\|$ is an eigenvalue they
converge to the corresponding eigenprojection.  What is asserted here is
only that the norm-level argument admits a counterexample.)

\smallskip\noindent\emph{The resolvent route stalls.}  A vector
$v \in \mathcal{D}_r(S)$ lies in the range of $t - S$ for every $t > r$ by
a Neumann series, but with a vector-dependent bound; and a self-adjoint
operator with dense range need not be invertible.

\smallskip\noindent The statement is about the spectral measure near the
top of the spectrum, and only a functional calculus sees that.

\begin{corollary}[support-exponential constants are not an obstruction]
\label{cor:noobstruction}
Let $\{A_s\}$ be observables whose span is dense in $H$, indexed so that
$A_s$ has support of size $s$, and suppose $\|S^n A_s\| \le
C(s)\,r^n$ for all $n$, with a rate $0 \le r < 1$ independent of $s$.  Then
$\|S\| \le r$, whatever the growth of $C$.  In particular $C(s) =
e^{\kappa s}$ is admissible.
\end{corollary}

Corollary~\ref{cor:noobstruction} is the reason the sharp form matters
for the programme this paper belongs to.  The volume-uniform area
law of~\cite{ErikssonAreaLaw} carries a constant of the form
$N_c\,e^{|\mathrm{supp}| \cdot 4d\,K}$, exponential in the edge support
of the loop, while its decay \emph{rate} is the support-independent
polymer activity rate; and the connecting-cluster estimate of the same
programme has a support-independent constant but covers only a local,
non-total family.  A natural reading of that pair is that the two horns
cannot be closed simultaneously, so that a total family must carry
growing constants and the gap must fail volume-uniformly.
Corollary~\ref{cor:noobstruction} shows that reading is false: the growth
of the constant is irrelevant, and only the uniformity of the
\emph{rate} is load-bearing.

We stress the limited role of that citation.  The area-law theorem is used
here \emph{only} to motivate the shape of an observable-dependent
constant.  It does not supply a dense family in any Osterwalder--Seiler
space, it does not supply the time-direction correlator, it does not
supply a common rate on such a family, and it does not supply the
identification with $T^n$.  Nothing in this paper brings a mass gap closer
to being deduced from an area law.

\section{The bridge in its own objects, and volume-uniformity}
\label{sec:vu}

Theorem~\ref{thm:sharp} is about a bare self-adjoint operator.  Composing
it with Definitions~\ref{def:data}--\ref{def:conn} gives the bridge in
the form an application would consume.

\begin{theorem}[sharp bridge]
\label{thm:sharpbridge}
Let $(T,\Om)$ be transfer data on a complex Hilbert space $H \neq 0$ with
$T$ self-adjoint, and let $r \ge 0$.  Then $\|S\| \le r$ if and only if
there is a set $D \subseteq H$ with dense span such that every $v \in D$
admits some finite $C$ with $|G_v(n)| \le C\,r^n$ for all $n$.
\end{theorem}

\begin{theorem}[volume-uniformity]
\label{thm:vu}
Let $(T_i,\Om_i)$ be transfer data on complex Hilbert spaces $H_i \neq 0$,
$i$ ranging over any index set, and let $0 < r < 1$ be a \emph{common}
rate for which each $i$ admits a family as in
Theorem~\ref{thm:sharpbridge} --- with per-observable constants that need
not be related across $i$ or within it.  Then, with $m = -\log r > 0$,
\[
  \|T_i - \ket{\Om_i}\bra{\Om_i}\| \;\le\; e^{-m}
  \qquad\text{for every } i .
\]
\end{theorem}

No uniform control of the constants $C_{v,i}$ is assumed.  The only
cross-volume quantitative input is the common rate $r$; once that rate is
fixed, Theorem~\ref{thm:sharpbridge} yields the same operator-norm bound
for every $i$, its conclusion being a bound whose constant never sees the
index.  In the intended reading $i$ is a finite volume and $m$ is a
volume-uniform mass; we stress that no transfer operator for any gauge
theory is constructed here, so the reading is an intention and not a
result.

\section{Non-vacuity}

Every statement above is exhibited with inhabitants in which the
fluctuation sector is nonzero, so that none of them is satisfied only in
the degenerate one-dimensional case.

\begin{proposition}
\label{prop:nonvac}
Let $\Om, w$ be orthonormal in $H$ and $0 < r < 1$.  Then
$T = r\,\mathrm{id} + (1-r)P$ is transfer data with
\[
  \|T - \ket{\Om}\bra{\Om}\| = r \quad\text{exactly},\qquad
  T - \ket{\Om}\bra{\Om} \neq 0,\qquad -\log r > 0 .
\]
The hypotheses are discharged concretely in $\R^2$ and in $\C^2$ with
$\Om = e_0$, $w = e_1$, so the family is inhabited.
\end{proposition}

\section{What remains}

Within the abstract implication isolated here, the remaining bridge
obligations are: reflection positivity of the Wilson gauge measure in the
time direction; the GNS quotient and the transfer operator with
$0 \le T \le 1$ and $T\Om = \Om$; the identification of the Euclidean
correlator with a matrix element; and density, in the fluctuation sector
of the resulting space, of a family of gauge-invariant observables for
which a clustering estimate at a common rate is available.  Each is
classical mathematics~\cite{OsterwalderSeiler,GlimmJaffe} that has not, to
our knowledge, been mechanized; none is established here.

We do not claim that list is exhaustive for any particular notion of
``mass gap''.  Depending on the formulation one may additionally need
uniqueness of the vacuum, control of the rate at all times rather than
along a subsequence, the correct gauge-invariant sector, the relation to
a Hamiltonian, and a volume or lattice-spacing limit.  What we do claim
is that the criterion proved above is one obligation that is now
discharged, and that the prefactor question is not among the remaining
ones.

\begin{remark}[a boundary of the formalization, not of the mathematics]
\label{rem:RC}
Theorem~\ref{thm:equiv} is formalized over a real Hilbert space and
Theorems~\ref{thm:sharp}--\ref{thm:vu} over a complex one.  The reason is
that the library's continuous functional calculus for self-adjoint
elements is derived from the complex case and is unavailable for
operators on a real space, and the library contains no complexification
of a real inner-product space.  The two are separate instantiations of
one argument; the transport between them is not formalized here and is
not claimed.  This is a boundary of the formalization, not of the
mathematics.
\end{remark}

\section{Formalization architecture}
\label{sec:arch}

The development is three modules and runs in one direction:
\begin{gather*}
\text{projection identities}
\;\to\;
\text{elementary equivalence}
\;\to\;
\text{decay subspace}\\
\to\;
\text{continuous cut-off}
\;\to\;
\text{sharp bridge.}
\end{gather*}

\paragraph{What the library did not provide.}  Two absences shaped the
design.  Mathlib has no complexification of a real inner-product space,
and its continuous functional calculus for self-adjoint elements is
derived from the complex case, so it is unavailable for operators on a
real Hilbert space.  Theorem~\ref{thm:equiv} needs no functional
calculus and is therefore stated over $\R$; Theorems~\ref{thm:sharp}
and~\ref{thm:sharpbridge} need one and are stated over $\C$.  That split
is forced by tooling, not by mathematics, and the transport between the
two is not formalized (Remark~\ref{rem:RC}).

\paragraph{Two devices that may be reusable.}  First, the $C^*$ step
$\|S^2\| = \|S\|^2$ for symmetric $S$ is proved directly from
Cauchy--Schwarz (\lean{norm\_sq\_of\_symm}), so the real-case module
depends on no $C^*$-algebra API at all.  Second, the cut-off quotient is
defined as
$k_n(x) = f(x)/(\max(c,|x|))^{2n}$ rather than $f(x)/x^{2n}$: the
denominator is bounded below by $c^{2n} > 0$ at \emph{every} real $x$, so
continuity needs no case analysis and no division by zero can occur even
when $0$ lies in the spectrum, while the even power makes the
factorisation $k_n(x)\,x^{2n} = f(x)$ hold identically, with no side
condition.  Both are small, and both removed what would otherwise have
been the fragile parts of the argument.

\paragraph{What verification actually costs.}  Building the three modules
and their closure takes $8160$ jobs; the full core root takes $8415$.  The
difference is not the interesting quantity: because the modules
\lean{import Mathlib} wholesale, essentially the entire cost in both cases
is Mathlib, retrieved from cache.  The lane's own content is three module
compilations.  We record this because a ``minimal verification target''
would be misleading here --- it would not meaningfully reduce the work a
referee must do, and the $8415$ figure should be read as evidence of
integration into the core, not as a measure of this paper's size.

\section{Reproducibility, verification links, and provenance}
\label{sec:prov}

\subsection*{Freeze}

All statements are machine-checked in Lean~4 (toolchain
\lean{leanprover/lean4:v4.29.0-rc6}) against Mathlib pinned to commit
\lean{0764272048\-0157414db592fa85b626dafb71355b}.  The freeze is the
source tree at commit
\lean{\anchor}; \lean{lake build YangMillsCore} completes with
\textbf{8415} jobs.  Every declaration reachable from the core root is
subject to an axiom oracle; the transcripts below record, for
\textbf{2304} \lean{\#print axioms} invocations, that the axiom sets
occurring are exactly \lean{[propext, Classical.choice, Quot.sound]},
\lean{[propext, Quot.sound]} and \lean{[propext]} --- all subsets of the
three permitted --- with zero occurrences of \lean{sorryAx} and no
project axioms.

\subsection*{Theorem-to-artifact map}

The three modules are
\osfile{TransferGap}{YangMills/OS/TransferGap.lean} (real case),
\osfile{DenseClustering}{YangMills/OS/DenseClustering.lean} (sharp form)
and \osfile{SharpBridge}{YangMills/OS/SharpBridge.lean} (sharp bridge);
the column below names the module by initial.

\small
\begin{longtable}{@{}>{\raggedright\arraybackslash}p{0.30\textwidth} >{\raggedright\arraybackslash}p{0.52\textwidth} l@{}}
\toprule
Statement here & Lean name & Mod. \\
\midrule
\endhead
Def.~\ref{def:data} ($P$, $S$) & \lean{vacuumProjection},
  \lean{projectedTransfer} & T \\
Def.~\ref{def:conn} ($G_v$) & \lean{connCorr} & T \\
transfer data & \lean{VacuumTransfer} & T \\
Lem.~\ref{lem:struct} & \lean{projected\_pow\_succ},
  \lean{connCorr\_eq} & T \\
Thm.~\ref{thm:equiv} & \lean{clustering\_iff\_gap} & T \\
$\|S^2\|=\|S\|^2$, proved by hand & \lean{norm\_sq\_of\_symm} & T \\
elementary dense form & \lean{gap\_of\_dense\_clustering} & T \\
Def.~\ref{def:decay} & \lean{decayDomain} & D \\
Lem.~\ref{lem:submodule} & \lean{decaySubmodule} & D \\
cut-off and factorisation & \lean{cutOff}, \lean{cutOffQuot},
  \lean{cutOffQuot\_mul\_pow}, \lean{cutOffQuot\_le} & D \\
Thm.~\ref{thm:sharp} & \lean{norm\_le\_of\_dense\_decayDomain},
  \lean{norm\_le\_of\_span\_dense\_decay} & D \\
Thm.~\ref{thm:sharpbridge} & \lean{sharp\_clustering\_iff\_gap} & S \\
Thm.~\ref{thm:vu} & \lean{volumeUniform\_sharp\_gap} (and
  \lean{volumeUniform\_gap}, real case) & S, T \\
Prop.~\ref{prop:nonvac} & \lean{exists\_vacuumTransfer\_gap},
  \lean{euclidean\_two\_dim\_vacuumTransfer\_gap} & T \\
Prop.~\ref{prop:nonvac} (complex) & \lean{exists\_vacuumTransferC\_gap},
  \lean{euclidean\_two\_dim\_vacuumTransferC\_gap} & S \\
\bottomrule
\end{longtable}
\normalsize

\subsection*{Oracle transcripts and build root}

\begin{itemize}
\item Core root:
  \repofile{YangMillsCore.lean}{YangMillsCore.lean}.
\item Oracle driver:
  \repofile{oracle\_check.lean}{oracle\_check.lean}.
\item Transcript at the real-case commit:
  \repofile{docs/oracle-transcripts/ORACLE-20260727-34696985.txt}{ORACLE-20260727-34696985.txt}.
\item Transcript at the sharp-form commit:
  \repofile{docs/oracle-transcripts/ORACLE-20260727-a0ce50ea.txt}{ORACLE-20260727-a0ce50ea.txt}.
\item Transcript at the sharp-bridge commit:
  \repofile{docs/oracle-transcripts/ORACLE-20260727-e5d76dae.txt}{ORACLE-20260727-e5d76dae.txt}.
\item Governance and scope record:
  \repofile{docs/O-BRIDGE-CHARTER.md}{docs/O-BRIDGE-CHARTER.md} and
  \repofile{docs/O-BRIDGE-AUDIT-20260727.md}{docs/O-BRIDGE-AUDIT-20260727.md}.
\end{itemize}

The transcripts were produced at the three code commits after which they
are named; those commits were replayed over two intervening
documentation-only commits before publication, and the replay is recorded
together with the byte-level equality of the three module blobs.

\subsection*{Provenance and prior art}

The equivalence between exponential clustering and a gap of the transfer
operator is \emph{classical} and no novelty is claimed for it: it is the
standard argument of constructive field
theory~\cite{GlimmJaffe,Simon}, and the lattice-gauge transfer operator
together with its positivity is due to Osterwalder and
Seiler~\cite{OsterwalderSeiler}.

\textbf{We make no claim of mathematical novelty for
Theorem~\ref{thm:sharp} either.}  Its content --- that
conceptually, $\mathcal{D}_r(S)$ is the spectral subspace
$\operatorname{ran} E([-r,r])$ of a self-adjoint $S$, so that its density
forces $\sigma(S) \subseteq [-r,r]$.  \emph{That identification is offered
as the conceptual reading and is not separately formalized here}; what the
Lean development contains is the norm criterion itself, as the
theorem-to-artifact map above records.  Criteria of this kind are standard
material in local spectral theory~\cite{LaursenNeumann}, where
$\limsup_n \|S^n v\|^{1/n}$ is the local spectral radius at $v$ and the
sets on which it is bounded are the local spectral subspaces.  We have not
located the precise statement in the form used here and have not performed
an exhaustive search; rather than assert priority we state the expectation
that it is known, and we invite the closest antecedent.  What we do claim
is the mechanization, the packaging of the criterion for \emph{families}
of observables with per-member constants, and
Corollary~\ref{cor:noobstruction}, which is the consequence that matters
for cluster expansions and which we have not seen drawn.

The volume-uniform area law whose constant shape motivates
Corollary~\ref{cor:noobstruction} is \cite{ErikssonAreaLaw}, the formal
antecedent within the same programme; its role here is motivational only,
as stated in Section~\ref{sec:sharp}.

\begin{thebibliography}{9}

\bibitem{OsterwalderSeiler}
K.~Osterwalder and E.~Seiler,
\emph{Gauge field theories on a lattice},
Ann. Physics \textbf{110} (1978), 440--471.

\bibitem{GlimmJaffe}
J.~Glimm and A.~Jaffe,
\emph{Quantum Physics: A Functional Integral Point of View},
2nd ed., Springer, 1987; in particular \S19.

\bibitem{Simon}
B.~Simon,
\emph{The $P(\phi)_2$ Euclidean (Quantum) Field Theory},
Princeton University Press, 1974.

\bibitem{LaursenNeumann}
K.~B.~Laursen and M.~M.~Neumann,
\emph{An Introduction to Local Spectral Theory},
London Mathematical Society Monographs New Series \textbf{20},
Clarendon Press / Oxford University Press, 2000.

\bibitem{ErikssonAreaLaw}
L.~Eriksson,
\emph{A volume-uniform area law for lattice Yang--Mills, machine-checked},
ai.viXra.org:2607.0005.

\bibitem{Lean4}
L.~de~Moura and S.~Ullrich,
\emph{The Lean 4 theorem prover and programming language},
CADE~28 (2021), 625--635.

\bibitem{Mathlib}
The mathlib Community,
\emph{The Lean mathematical library},
CPP 2020, 367--381.

\bibitem{Clay}
A.~Jaffe and E.~Witten,
\emph{Quantum Yang--Mills theory},
Clay Mathematics Institute Millennium Problem description, 2000.

\end{thebibliography}

\end{document}
